\documentclass[12pt]{article}
\usepackage[T2A]{fontenc}
\usepackage[utf8]{inputenc}
\usepackage[russian]{babel}
\usepackage{latexsym,amssymb,amsthm}
\usepackage{amsmath}
\usepackage{wrapfig}
\RequirePackage[pdftex]{graphicx}
%%\usepackage{array}

\usepackage{epsf,epic,eepic}
% \input ris.tex
% \input cells.sty
% \addrisoption{\risleftfalse}

\DeclareRobustCommand{\No}{\ifmmode{\nfss@text{\textnumero}}\else\textnumero\fi}

\oddsidemargin=0pt
\textwidth=159mm
\headheight=0pt
\headsep=0pt
\topmargin=0pt
\textheight=240mm

%%
%% Три точки для обозначения делимости
%%
\catcode`\@=11
\def\kratno{\mathrel{\smash{\lower.5ex\hbox{$\,\vdots\,$}}}}
% Это прямые три точки для обозначения делимости
\def\nekratno{\mathrel{\mathpalette\c@ncel\kratno}}
\def\c@ncel#1#2{\m@th\ooalign{$\hfil#1\mkern1mu/\hfil$\crcr$#1#2$}}
% Это перечеркнутые три точки для обозначения неделимости
\def\divis{\smash{\lower.1ex\hbox{\,\hbox{.}\kern-.22em\lower-.6ex\hbox{.}
\kern-.52em\lower-1.2ex\hbox{.}\,}} }
% А это наклонные...

\hyphenation{чис-ло чис-ла чис-лу чис-лом чис-ле чис-лам чис-ла-ми чис-лах
че-ты-рех-уголь-ник че-ты-рех-уголь-ни-ка пря-мо-уголь-ник тре-уголь-ни-ка
тре-уголь-ни-ков вне-впи-сан-ной вне-впи-сан-ных}

\def\signed(#1){{\unskip\nobreak\hfil\penalty50\hskip1em\hbox{}\nobreak\hfil
                 \rm(\small\it#1\/\rm)\parfillskip=0pt\finalhyphendemerits=0\par}}
        % macro for authorship's signature, from TeXbook
\def\be{\signed(С.~Берлов)}
\def\fp{\signed(Ф.~Петров)}
\def\frank{\signed(В.~Франк)}
\def\bah{\signed(Ф.~Бахарев)}
\def\guas{\signed(А.~Голованов)}
\def\as{\signed(А.~Смирнов)}
\def\dvk{\signed(Д.~Карпов)}
\def\dm{\signed(Д.~Максимов)}
\def\kk{\signed(К.~Кохась)}
\def\kuz{\signed(А.~Кузнецов)}
\def\hr{\signed(А.~Храбров)}
\def\ov{\signed(О.~Ванюшина)}
\def\io{\signed(О.~Иванова)}
\def\sop{\signed(Е.~Сопкина)}
\def\sviv{\signed(С.~Иванов)}
\def\suk{\signed(К.~Сухов)}
\def\dr{\signed(Д.~Ростовский)}
\def\pas{\signed(А.~Пастор)}
\def\azh{\signed(А.~Железняк)}
\def\vlas{\signed(Н.~Власова)}
\def\asol{\signed(А.~Солынин)}
\def\miv{\signed(М.~Иванов)}
\def\chuh{\signed(А.~Чухнов)}
\def\lev{\signed(Л.~Емельянов)}
\def\dsh{\signed(Д.~Ширяев)}
\def\zh{\signed(жюри)}
\def\bdzh{\signed(О.~Бадажкова)}

\sloppy
\pagestyle{empty}

\def\q#1.{{\bf #1. }}
\def\dotline{\smallskip\hbox to \hsize{\dotfill}\medskip}

\begin{document}

\centerline{\sc САНКТ-ПЕТЕРБУРГСКАЯ}
\centerline{\sc ОЛИМПИАДА ШКОЛЬНИКОВ 2021 года ПО МАТЕМАТИКЕ}
\centerline{\sc II тур. \ 9 класс.}
\medskip
\hrule
\bigskip

\q1. На окружности поставлена 2021 точка.
Костя отмечает точку, потом отмечает соседнюю справа точку, 
потом он отмечает точку справа через одну от только что отмеченной, 
потом~--- точку справа через две от только что отмеченной и т.\,д.
На каком ходу впервые появится точка, отмеченная дважды? 
\kk
 
\q2. У Миши есть шахматная доска   $100\times 100$ и мешок со 199 ладьями.  За один ход можно либо поставить 
одну ладью из мешка на левую нижнюю клетку, либо снять две ладьи, стоящие на одной клетке, 
одну из них поставить на клетку, соседнюю сверху или справа, а другую убрать в мешок.  
Миша хочет расставить на доске 100 не бьющих друг друга ладей.
Любую ли такую расстановку он сможет получить?
\signed(М.~Иванов, С.~Берлов)  
 
\q3. Во вписанном четырехугольнике $ABCD$ \  ${\angle A = 3 \angle B}$. На стороне $AB$ отмечена точка $C_1,$ а на стороне $BC$~--- точка $A_1$ так, что $AA_1 = AC = CC_1$. Докажите, что $3A_1C_1 > BD$. 
\kuz

\q4. На плоскости отмечено $n$ точек с разными абсциссами. Через
каждую пару точек провели {\em параболу}~---  график квадратного трехчле\-на с~единичным старшим коэффициентом. Парабола называется
хорошей, если ни на ней, ни выше неё нет отмеченных точек, кроме
тех двух, через которые она проведена. Какое наибольшее количество хороших парабол могло получиться?
\signed(Е.~Кравченко )

\dotline

\centerline{Олимпиада 2021 года. II тур. 9 класс. Выводная аудитория.}
\smallskip

\q5. Равнобедренная трапеция $ABCD$, в которой основание $AD$ вдвое
больше боковой стороны $AB$, вписана в окружность $\omega$.
Точки $E$ и $F$ выбраны на окружности $\omega$ так, что $AC
\parallel DE$ и $BD \parallel AF$. Отрезок $BE$ пересекает
отрезки $AC$ и $AF$ в точках $X$ и $Y$ соответственно.
Докажите, что описанные окружности треугольников $BCX$ и
$EFY$ касаются. 
\kuz

\q6. В школе 450 учеников. Каждый из них имеет не меньше 100 друзей среди остальных, и среди любых 200 учеников всегда найдутся двое друзей. Докажите, что можно отправить в байдарочный поход 302 ученика так, чтобы в каждой из 151 двухместной байдарки оказались друзья. 
\dvk 
 
\q7. Дано натуральное число $n$. Для каждого простого числа $p$ из промежутка
$[n,n^4]$ посчитали число $\frac{1}{p}$, и все полученные числа сложили.
Докажите, что их сумма меньше $4$.
\signed(Н.~Филонов)


\clearpage

\centerline{\sc САНКТ-ПЕТЕРБУРГСКАЯ}
\centerline{\sc ОЛИМПИАДА ШКОЛЬНИКОВ 2021 года ПО МАТЕМАТИКЕ}
\centerline{\sc II тур. \ 10 класс.}
\medskip
\hrule
\medskip

\q1. Решите систему уравнений: 
\begin{align*}
\sin^2x+ \cos^2y&=y^4, \\ 
\sin^2y+ \cos^2x&=x^2.
\end{align*}
\hr

\q2. В ряд выписано 2021 простое натуральное число. Каждое, кроме крайних, отличается от одного из своих соседей на 12, а от другого~--- на 6. Докажите, что среди этих чисел есть равные.
\signed(М.~Антипов)

\q3. Дан выпуклый пятиугольник $ABCDE$. Точки
$A_1$, $B_1$, $C_1$, $D_1$, $E_1$ таковы, что 
$$
AA_1\perp EB,\quad BB_1\perp AC,\quad CC_1\perp BD,\quad DD_1\perp CE,\quad EE_1\perp DA.
$$
Кроме того, $AE_1=AB_1$, $BC_1=BA_1$, $CB_1=CD_1$, $DC_1=DE_1$.  Докажите, что тогда $ED_1=EA_1$.

\q4. Штирлиц хочет послать в Центр шифровку, представляющую собой
код из 100 символов <<точка>> или <<тире>>.
Полученная им накануне из Центра \emph{Инструкция о
конспирации} гласит:

--- при передаче шифровки по радио ровно 49 символов следует
заменить на противоположные;

--- расположение <<неверных>> символов возлагается на
передающую сторону и с Центром не обсуждается.

Докажите, что Штирлиц может послать свою шифровку 10 раз,
подбирая при каждой передаче 49 символов так, чтобы Центр,
получив эти 10 шифровок, имел возможность однозначно
восстановить исходный код.

\dotline

\centerline{Олимпиада 2021 года. II тур. 10 класс. Выводная аудитория.}
\smallskip

\q5.  Вершины выпуклого 2550-угольника покрашены в черный и белый цвета так: чёрная, белая, две чёрные, две белые, три чёрные, три белые, $\dots$, 50 чёрных, 50 белых. Даня разрезал его на четырёхугольники диагоналями, не
имеющими общих внутренних точек. Докажите, что найдётся четырёхугольник разрезания, в котором две соседние вершины чёрные, а две другие вершины~--- белые.  
\signed(Д.~Руденко)

\q6. Прямая $\ell$ проходит через вершину $C$ ромба $ABCD$ и
пересекает продолжения его сторон $AB$ и $AD$ в точках $X$ и
$Y$ соответственно. Прямые $DX$ и $BY$ вторично пересекают
окружность, описанную около треугольника $AXY$, в точках $P$
и $Q$. Докажите, что окружность,  описанная около
треугольника $PCQ$, касается прямой~$\ell$.
\kuz

\q7. Коля нашёл несколько попарно взаимно простых натуральных чисел, каждое из которых меньше квадрата любого другого. Докажите, что сумма величин, обратных к Колиным числам, меньше~2.


\clearpage

\centerline{\sc САНКТ-ПЕТЕРБУРГСКАЯ}
\centerline{\sc ОЛИМПИАДА ШКОЛЬНИКОВ 2021 года ПО МАТЕМАТИКЕ}
\centerline{\sc II тур. \ 11 класс.}
\medskip
\hrule
\bigskip

\q1. Дано простое число $p$. Все натуральные числа от $1$ до $p$ выписаны в ряд в порядке возрастания. 
Найдите все $p$, для которых этот ряд можно разбить на несколько блоков подряд идущих чисел так, чтобы суммы чисел во всех блоках были одинаковы. 
\hr

\q2. Клетки таблицы $100\times 100$ окрашены в белый цвет.
За один ход разрешается выбрать любые $99$ клеток из одной строки или из
одного столбца и~перекрасить каждую из них в противоположный цвет~--- из белого в черный, а из черного в
белый. За какое наименьшее количество ходов можно получить таблицу с шахматной раскраской клеток?
\be

\q3.
В пирамиде $SA_1A_2\dots A_n$ все боковые рёбра равны. Точка~$X_1$~--- середина дуги $A_1A_2$ описанной окружности треугольника $SA_1A_2$, точка~$X_2$~--- середина дуги $A_2A_3$ описанной окружности треугольника
$SA_2A_3$ и~т.\,д., точка $X_n$~--- середина дуги $A_nA_1$ описанной окружности треугольника 
$SA_nA_1$. Докажите, что описанные ок\-руж\-ности треугольников \ $X_1A_2X_2$, \ $X_2A_3X_3$, \ \dots, \ $X_nA_1X_1$ пересекаются в одной точке.

\q4.
На доске написаны функции 
$$
F(x)=x^2+\frac{12}{x^2},\qquad G(x)=\sin(\pi x^2) \qquad\text{и}\qquad H(x)=1.
$$ 
Если на доске уже написаны функции $f(x)$ и $g(x)$, то можно выписать на доску еще и функции 
$$
f(x)+g(x), \qquad f(x)-g(x), \qquad f(x)g(x), \qquad cf(x)
$$
(последнюю~--- с любым вещественным коэффициентом $c$). 
Может ли на доске появиться такая функция $h(x)$, что $|h(x)-x|<\frac{1}{3}$ при всех ${x\in[1,10]}$?

\dotline

\centerline{Олимпиада 2021 года. II тур. 11 класс. Выводная аудитория.}
\smallskip

\q5. Дано натуральное число $n$. Для каждого простого числа $p$ из промежутка
$[n,n^2]$ посчитали число $\frac{1}{p}$, и все полученные числа сложили.
Докажите, их сумма меньше $2^{\vphantom{2}}$.

\q6. Точка $M$~--- середина основания $AD$ трапеции $ABCD$, вписанной в
ок\-руж\-ность $\omega$. Биссектриса угла $ABD$ пересекает отрезок~$AM$ в
точке~$K$. Прямая $CM$ вторично пересекает окружность~$\omega$ в~точке
$N$. Из точки $B$ проведены касательные $BP$ и $BQ$ к описанной окружности
треугольника $MKN$. Докажите, что прямые $BK$, $MN$ и $PQ$ пересекаются в~одной точке.
\kuz

\q7. Квадрат разрезан на красные и синие прямоугольники. Сумма площадей красных
прямоугольников равна сумме площадей синих. Для каждого синего
прямоугольника запишем отношение длины его вертикальной стороны к~длине
горизонтальной, а для каждого красного  прямоугольника~--- отношение 
длины его горизонтальной стороны к длине вертикальной. Найдите наименьшее возможное 
значение суммы всех записанных чисел.


\end{document}
